% GENERATED FILE. Edit canonical lessons, then run npm run generate:books.
% canonical-source-hash: fnv1a64:e0986c08ccc49bd6
% canonical-lessons: HI-C503-grhasthi, HI-C503-mukhiya, HI-C503-buzurg, HI-C503-pirhi, HI-C503-posna

\chapter{The Generations}
\label{ch:hi-the-generations}
\hlchaptermodality{503}

\begin{chapteropening}
\textbf{By the end of this chapter:} \emph{I can talk about the generations of a family.}

Complete the last lesson of chapter 503: I can talk about the generations of a family.
\end{chapteropening}

\section[\texorpdfstring{\hi{गृहस्थी} (gṛhasthī)}{गृहस्थी (gṛhasthī)}]{\hi{गृहस्थी} (gṛhasthī) — a household; family life}

\label{lesson:HI-C503-grhasthi}



\begin{quote}
Before the new one: say the Hindi for to rock, then the Hindi for only (child).
\end{quote}

\subsection*{You'll want to know: \hi{गृहस्थी}}
\textbf{\hi{गृहस्थी}} — \emph{gṛhasthī} — \textquotedblleft{}a household; family life\textquotedblright{}.

\begin{quote}
\hi{शादी} \hi{के} \hi{बाद} \hi{मीरा} \hi{ने} \hi{अपनी} \hi{गृहस्थी} \hi{सँभाली।}
\end{quote}

\begin{quote}
\emph{śādī ke bād Mīrā ne apnī gṛhasthī sãbhālī.}
\end{quote}

\begin{quote}
After her marriage Mira took charge of her household.
\end{quote}

Feminine.

\begin{cousinweb}
\textbf{\hi{गृहस्थी}} (\emph{gṛhasthī}) is Sanskrit \emph{gṛhastha}, \textquotedblleft{}a householder\textquotedblright{}, built on \emph{gṛha}, \textquotedblleft{}house\textquotedblright{}, and \emph{stha}, \textquotedblleft{}standing in\textquotedblright{} (Monier-Williams), with the ending \emph{-ī}: the life and business of a household.
\end{cousinweb}

\begin{grammarlens}[title={What you've built}]
One more word for describing what you lived and felt.
\end{grammarlens}

\subsection*{Your turn}
\begin{itemize}
\raggedright
  \item \emph{Say it:} \emph{gṛhasthī}
  \item \emph{Say it:} \emph{gṛhasthī}, once more
  \item \emph{Say it:} \emph{śādī ke bād Mīrā ne apnī gṛhasthī sãbhālī.} — After her marriage Mira took charge of her household.
  \item \emph{Recall:} say the Hindi for to rock, then the Hindi for only (child), then say \emph{gṛhasthī} again
\end{itemize}

\begin{tcolorbox}[breakable,colback=teal!4,colframe=teal!35!black,title={Before you move on}]
What does \hi{गृहस्थी} mean? (\textquotedblleft{}A household; family life\textquotedblright{}.) Say it once more.
\end{tcolorbox}
\section[\texorpdfstring{\hi{मुखिया} (mukhiyā)}{मुखिया (mukhiyā)}]{\hi{मुखिया} (mukhiyā) — the head (of a family or village)}

\label{lesson:HI-C503-mukhiya}



\begin{quote}
Before the new one: say the Hindi for only (child), then the Hindi for a household.
\end{quote}

\subsection*{You'll want to know: \hi{मुखिया}}
\textbf{\hi{मुखिया}} — \emph{mukhiyā} — \textquotedblleft{}the head (of a family or village)\textquotedblright{}.

\begin{quote}
\hi{दादा} \hi{हमारे} \hi{परिवार} \hi{के} \hi{मुखिया} \hi{हैं।}
\end{quote}

\begin{quote}
\emph{dādā hamāre parivār ke mukhiyā haiṁ.}
\end{quote}

\begin{quote}
Grandfather is the head of our family.
\end{quote}

Masculine.

\begin{cousinweb}
\textbf{\hi{मुखिया}} (\emph{mukhiyā}) is a tadbhava: CDIAL lists Hindi \emph{mukhyā}, \textquotedblleft{}chief\textquotedblright{}, under Sanskrit \emph{múkhya}, \textquotedblleft{}pertaining to the face\textquotedblright{} and \textquotedblleft{}chief\textquotedblright{} (CDIAL 10174); \emph{mukhiyā} is the usual spelling. \emph{múkhya} is built on \emph{múkha}, \textquotedblleft{}face, front\textquotedblright{}, the word behind \emph{mũh}.
\end{cousinweb}

\begin{grammarlens}[title={What you've built}]
One more word for describing what you lived and felt.
\end{grammarlens}

\subsection*{Your turn}
\begin{itemize}
\raggedright
  \item \emph{Say it:} \emph{mukhiyā}
  \item \emph{Say it:} \emph{mukhiyā}, once more
  \item \emph{Say it:} \emph{dādā hamāre parivār ke mukhiyā haiṁ.} — Grandfather is the head of our family.
  \item \emph{Recall:} say the Hindi for only (child), then the Hindi for a household, then say \emph{mukhiyā} again
\end{itemize}

\begin{tcolorbox}[breakable,colback=teal!4,colframe=teal!35!black,title={Before you move on}]
What does \hi{मुखिया} mean? (\textquotedblleft{}The head (of a family or village)\textquotedblright{}.) Say it once more.
\end{tcolorbox}
\section[\texorpdfstring{\hi{बुज़ुर्ग} (buzurg)}{बुज़ुर्ग (buzurg)}]{\hi{बुज़ुर्ग} (buzurg) — an elder; elderly}

\label{lesson:HI-C503-buzurg}



\begin{quote}
Before the new one: say the Hindi for a household, then the Hindi for the head (of a family or village).
\end{quote}

\subsection*{You'll want to know: \hi{बुज़ुर्ग}}
\textbf{\hi{बुज़ुर्ग}} — \emph{buzurg} — \textquotedblleft{}an elder; elderly\textquotedblright{}.

\begin{quote}
\hi{गाँव} \hi{के} \hi{बुज़ुर्ग} \hi{पेड़} \hi{के} \hi{नीचे} \hi{बैठकर} \hi{बात} \hi{करते} \hi{हैं।}
\end{quote}

\begin{quote}
\emph{gā̃v ke buzurg peṛ ke nīce baiṭhkar bāt karte haiṁ.}
\end{quote}

\begin{quote}
The village elders sit under the tree and talk.
\end{quote}

Masculine or feminine, by the person meant.

\begin{cousinweb}
\textbf{\hi{बुज़ुर्ग}} (\emph{buzurg}) is usually traced to Persian \emph{buzurg}, \textquotedblleft{}great\textquotedblright{}, and so \textquotedblleft{}an elder\textquotedblright{}. It is a respectful word for anyone old.
\end{cousinweb}

\begin{grammarlens}[title={What you've built}]
One more word for describing what you lived and felt.
\end{grammarlens}

\subsection*{Your turn}
\begin{itemize}
\raggedright
  \item \emph{Say it:} \emph{buzurg}
  \item \emph{Say it:} \emph{buzurg}, once more
  \item \emph{Say it:} \emph{gā̃v ke buzurg peṛ ke nīce baiṭhkar bāt karte haiṁ.} — The village elders sit under the tree and talk.
  \item \emph{Recall:} say the Hindi for a household, then the Hindi for the head (of a family or village), then say \emph{buzurg} again
\end{itemize}

\begin{tcolorbox}[breakable,colback=teal!4,colframe=teal!35!black,title={Before you move on}]
What does \hi{बुज़ुर्ग} mean? (\textquotedblleft{}An elder; elderly\textquotedblright{}.) Say it once more.
\end{tcolorbox}
\section[\texorpdfstring{\hi{पीढ़ी} (pīṛhī)}{पीढ़ी (pīṛhī)}]{\hi{पीढ़ी} (pīṛhī) — a generation}

\label{lesson:HI-C503-pirhi}



\begin{quote}
Before the new one: say the Hindi for the head (of a family or village), then the Hindi for an elder.
\end{quote}

\subsection*{You'll want to know: \hi{पीढ़ी}}
\textbf{\hi{पीढ़ी}} — \emph{pīṛhī} — \textquotedblleft{}a generation\textquotedblright{}.

\begin{quote}
\hi{हमारे} \hi{घर} \hi{में} \hi{तीन} \hi{पीढ़ियाँ} \hi{एक} \hi{साथ} \hi{रहती} \hi{हैं।}
\end{quote}

\begin{quote}
\emph{hamāre ghar meṁ tīn pīṛhiyā̃ ek sāth rahtī haiṁ.}
\end{quote}

\begin{quote}
Three generations live together in our house.
\end{quote}

Feminine.

\begin{cousinweb}
\textbf{\hi{पीढ़ी}} (\emph{pīṛhī}) is a tadbhava: CDIAL lists Hindi \emph{pīṛhī} both as \textquotedblleft{}small square stool\textquotedblright{} and as \textquotedblleft{}generation\textquotedblright{} under Sanskrit \emph{pīṭha}, \textquotedblleft{}seat\textquotedblright{} (CDIAL 8222).
\end{cousinweb}

\begin{grammarlens}[title={What you've built}]
One more word for describing what you lived and felt.
\end{grammarlens}

\subsection*{Your turn}
\begin{itemize}
\raggedright
  \item \emph{Say it:} \emph{pīṛhī}
  \item \emph{Say it:} \emph{pīṛhī}, once more
  \item \emph{Say it:} \emph{hamāre ghar meṁ tīn pīṛhiyā̃ ek sāth rahtī haiṁ.} — Three generations live together in our house.
  \item \emph{Recall:} say the Hindi for the head (of a family or village), then the Hindi for an elder, then say \emph{pīṛhī} again
\end{itemize}

\begin{tcolorbox}[breakable,colback=teal!4,colframe=teal!35!black,title={Before you move on}]
What does \hi{पीढ़ी} mean? (\textquotedblleft{}A generation\textquotedblright{}.) Say it once more.
\end{tcolorbox}
\section[\texorpdfstring{\hi{पोसना} (posnā)}{पोसना (posnā)}]{\hi{पोसना} (posnā) — to nurture; to rear}

\label{lesson:HI-C503-posna}



\begin{quote}
Before the new one: say the Hindi for an elder, then the Hindi for a generation.
\end{quote}

\subsection*{You'll want to know: \hi{पोसना}}
\textbf{\hi{पोसना}} — \emph{posnā} — \textquotedblleft{}to nurture; to rear\textquotedblright{}.

\begin{quote}
\hi{दादी} \hi{ने} \hi{हमारे} \hi{परिवार} \hi{की} \hi{तीन} \hi{पीढ़ियों} \hi{को} \hi{पोसा।}
\end{quote}

\begin{quote}
\emph{dādī ne hamāre parivār kī tīn pīṛhiyoṁ ko posā.}
\end{quote}

\begin{quote}
Grandmother has nurtured three generations of our family.
\end{quote}

Past, with \hi{ने} (he or she …): \hi{उसने} \hi{बच्चों} \hi{को} \hi{पोसा}, \emph{usne baccoṁ ko posā}.

\begin{cousinweb}
\textbf{\hi{पोसना}} (\emph{posnā}) is a tadbhava: CDIAL lists Hindi \emph{posnā}, \textquotedblleft{}to nourish, rear\textquotedblright{}, under Sanskrit \emph{pōṣáyati}, \textquotedblleft{}nourishes\textquotedblright{} (CDIAL 8410).
\end{cousinweb}

\begin{grammarlens}[title={What you've built}]
One more word for describing what you lived and felt.
\end{grammarlens}

\subsection*{Your turn}
\begin{itemize}
\raggedright
  \item \emph{Say it:} \emph{posnā}
  \item \emph{Say it:} \emph{posnā}, once more
  \item \emph{Say it:} \emph{dādī ne hamāre parivār kī tīn pīṛhiyoṁ ko posā.} — Grandmother has nurtured three generations of our family.
  \item \emph{Recall:} say the Hindi for an elder, then the Hindi for a generation, then say \emph{posnā} again
\end{itemize}

\begin{tcolorbox}[breakable,colback=teal!4,colframe=teal!35!black,title={Before you move on}]
What does \hi{पोसना} mean? (\textquotedblleft{}To nurture; to rear\textquotedblright{}.) Say it once more.
\end{tcolorbox}
